\documentclass[12pt]{article}
\usepackage[utf8]{inputenc}
\usepackage[T1]{fontenc}
\usepackage{hyperref}
\title{Classical Tibetan Language Resources for K-12 Curriculum Development}
\author{Compiled by Claude Sonnet 5.5 (claude-sonnet-5-5) for Andr\'as Kornai}
\date{}
\begin{document}
\maketitle

\section{Introduction}
Classical Tibetan (ISO 639-3 xct; Glottocode clas1254) is the literary language of Tibetan Buddhist and other texts. It has no native speakers but a very large body of literature. Search turned up many university-level textbooks and digital libraries, but nothing designed for K-12 schools.

\section{Primary Educational Resources}
\begin{itemize}
\item Hackett, \emph{Learning Classical Tibetan} (Snow Lion): classroom and self-study textbook. Each chapter pairs a Buddhist text with grammar notes, vocabulary and a literal translation. \url{https://www.shambhala.com/learning-classical-tibetan.html}
\item Bialek, \emph{A Textbook in Classical Tibetan} (Routledge, 2022), and Wilson, \emph{Translating Buddhism from Tibetan}, are on the University of Edinburgh's course reading list. I did not see a publisher URL for either, only the course page. \url{https://www.drps.ed.ac.uk/26-27/dpt/cxdivi10126.htm}
\item A 2011 Indology mailing-list thread recommending Wilson and Hahn's \emph{Textbook of Classical Literary Tibetan}. The advice is dated. \url{https://list.indology.info/pipermail/indology/2011-November/036021.html}
\end{itemize}

\section{Linguistic and Reference Materials}
\begin{itemize}
\item Monlam Online Dictionary, listed as a reference on the Edinburgh course. \url{https://monlamdictionary.com/}
\item Old Tibetan Documents Online (OTDO), ILCAA, Tokyo University of Foreign Studies: edited Old Tibetan texts (7th--12th c.) with KWIC search. This is older than Classical Tibetan. \url{https://otdo.aa-ken.jp}
\end{itemize}

\section{Religious and Community Resources}
\begin{itemize}
\item Buddhist Digital Resource Center (BDRC): a nonprofit that preserves and distributes Tibetan and other Buddhist texts. One source describes it as offering over 25 million pages, free to all. \url{https://en.wikipedia.org/wiki/Buddhist_Digital_Resource_Center}
\item Tsadra Foundation's Rangjung Yeshe Wiki resources page for translators, scholars and practitioners. \url{https://rywiki.tsadra.org/index.php/Resources}
\end{itemize}

\section{Digital and Audio Resources}
BDRC and OTDO are the digital corpora above. A search result mentioned an annotated corpus of Classical Tibetan (ACTib), built on BDRC texts and not manually corrected. I did not confirm its URL and do not cite one. I found no audio.

\section{Curriculum Development Context}
All materials found are for adult or university learners, and I found no K-12 curriculum. A school programme could use the textbooks above for graded readings. Classical Tibetan is also taught through the study of Buddhist texts, and BDRC offers a large source of authentic material.
\end{document}